\bookchapter{Bab 3: Menghubungkan Ponsel ke “Pancaran Kehidupan” — Jaringan}

\chapterintro{Tanpa jaringan, ponsel pintar setengah “tumpul”. Jaringan adalah garis kehidupan ponsel modern; memahaminya, menggunakannya dengan baik, dan melindunginya adalah pelajaran wajib bagi setiap pengguna ponsel.}

\booksection{Terhubung Tidak Bingung — Bedakan “Jalan” dan “Kendaraan” Wi-Fi dan Kuota Data}

Banyak orang melihat ikon Wi-Fi dan sinyal menyala bersamaan di bagian atas layar ponsel mereka, tetapi bingung tentang perbedaan esensialnya. Kesalahpahaman paling klasik adalah:\textbf{terhubung ke Wi-Fi tetapi tidak bisa membuka halaman}— karena router kehilangan koneksi internet, sementara ponsel tidak otomatis beralih ke kuota data. Untuk benar-benar memahami masalah ini, kita perlu membangun model kognitif “jalan dan kendaraan”.

\begin{tcbinfo}{Analogi Inti: Jalan dan Kendaraan}
Bayangkan internet sebagai sebuah kota.\textbf{“Jalan”}adalah jalur antara ponsel Anda dan internet,\textbf{“kendaraan”}adalah data yang berjalan di jalan tersebut. Wi-Fi dan kuota data adalah dua jalan yang berbeda, tetapi data (kendaraan) itu sendiri sama.

\end{tcbinfo}

\booksubsection{Perbedaan Esensial: Saluran Lokal vs Jalan Tol Jarak Jauh}

Wi-Fi pada dasarnya adalah\textbf{teknologi jaringan area lokal}, seperti jalan nirkabel yang dipasang di bawah gedung rumah Anda, bertanggung jawab mengirim data dari router ke ponsel. Namun jika lebar pita menunggak atau kabel serat optik putus, jalan masih ada, tetapi “jalan tol” di luar terputus, jadi Anda tidak bisa ke mana pun. Sedangkan kuota data (data seluler) ditransmisikan melalui menara pemancar 4G/5G operator, seperti langsung masuk ke jalur khusus operator, selama ada sinyal Anda bisa langsung terhubung ke internet.

\booksubsection{Biaya dan Skenario: Langganan Bulanan Tetap vs Biaya Berdasarkan Penggunaan}

Wi-Fi biasanya investasi sekali (biaya bulanan lebar pita), selama listrik tidak padam, seluruh keluarga bisa pakai tanpa batas, sangat cocok untuk skenario kuota data besar seperti mengunduh film atau bermain game. Sedangkan kuota data ditagih berdasarkan pemakaian, jika melebihi paket mungkin kecepatan diturunkan atau dikenakan biaya tambahan, lebih cocok untuk tugas ringan saat di kereta bawah tanah, perjalanan dinas, dan skenario mobile lainnya.

\begin{tcbpractice}{Mari Coba: Periksa Paket Kuota Data Anda}
Buka [Pengaturan] > [Jaringan Seluler] > [Manajemen Kuota Data] di ponsel, lihat berapa banyak kuota data tersisa bulan ini dan aplikasi mana yang paling banyak mengonsumsi kuota. Tahu posisinya, baru tidak panik saat keluar rumah.

\end{tcbpractice}

\booksubsection{Pertimbangan Keamanan: Wi-Fi Publik adalah “Kolam Pancing” Tersembunyi}

Dalam hal keamanan, data seluler biasanya lebih aman daripada Wi-Fi publik. Data seluler dienkripsi oleh operator, sedangkan Wi-Fi publik memang mudah dijadikan “hotspot palsu” oleh peretas untuk mencegat data yang tidak terenkripsi. Namun di sini perlu dibedakan: di era ketika HTTPS sudah meluas, yang dapat dicegat penyerang adalah\textbf{situs web mana yang Anda kunjungi}(domain), dan\textbf{konten HTTP yang tidak terenkripsi}; sedangkan pesan WeChat, kata sandi login aplikasi bank, dan data lain yang dienkripsi dengan HTTPS, penyerang\textbf{tidak dapat langsung membacanya}— kecuali melakukan serangan SSL stripping yang lebih kompleks dan pengguna mengabaikan peringatan sertifikat. Namun demikian, saat melakukan pembayaran di luar, menggunakan kuota data ponsel tetap pilihan paling aman.

\textbf{Peringatan Keamanan}\par

Di tempat umum seperti bandara, mal, kafe, jangan hubungkan ke Wi-Fi gratis dengan nama mencurigakan (seperti “Free\_WiFi\_2024”). Bahkan jika terhubung, jangan pernah login ke aplikasi bank atau memasukkan kata sandi di dalamnya.

\booksubsection{Penilaian di Sisi Ponsel: Pahami “Kode Rahasia” di Status Bar}

Bagaimana mengetahui apa yang sedang digunakan ponsel untuk terhubung? Mengambil ponsel Huawei sebagai contoh, munculnya\textbf{ikon panah dua arah}di samping sinyal operator atau sinyal WLAN di status bar menandakan sedang menggunakan metode tersebut untuk internet; jika keduanya menyala bersamaan, berarti akselerasi jaringan aktif, ponsel sedang memanfaatkan dua saluran sekaligus.

Pengguna iPhone lebih intuitif: jika ikon Wi-Fi menyala berarti menggunakan Wi-Fi, jika tidak menyala lihat batang sinyal di kanan atas berarti menggunakan kuota data. Sebagian ponsel Android juga akan menampilkan “: Jika sinyal buruk, secara manual beralih kembali ke 4G, daya tahan baterai akan jauh lebih baik.” atau “4G” di status bar, langsung memberi tahu jenis jaringan seluler saat ini.

\booksection{Jaringan Seluler (4G/5G) dan Kuota Data}

Jaringan seluler adalah dasar Anda terhubung kapan saja dan di mana saja. Dari 2G hingga 5G, kecepatan meningkat ribuan kali, tetapi prinsipnya sebenarnya mirip.

\booksubsection{Standar Jaringan:}

\begin{itemize}
  \item \textbf{4G LTE}: Saat ini paling luas cakupannya, kecepatan cukup (secara teoritis tertinggi 100Mbps+), konsumsi daya relatif rendah.
  \item \textbf{: Jika sinyal buruk, secara manual beralih kembali ke 4G, daya tahan baterai akan jauh lebih baik.}: Kecepatan tinggi (secara teoritis tertinggi 10Gbps), tetapi cakupan masih dalam penyempurnaan, konsumsi daya tinggi. Cocok untuk mengunduh file besar, menonton video HD.
  \item \textbf{5G SA vs NSA}: SA adalah jaringan mandiri, NSA bergantung pada menara 4G. Pengalaman SA lebih baik, tetapi membutuhkan dukungan operator dan ponsel.
\end{itemize}

\booksubsection{Manajemen Kuota Data:}

\begin{itemize}
  \item \textbf{Apa itu kuota data}: Yaitu jumlah data yang dihasilkan saat ponsel terhubung ke internet. Menonton video dan men-scroll gambar paling boros kuota, teks dan suara relatif hemat.
  \item \textbf{Tips hemat kuota}: Aktifkan "Mode Hemat Data" di pengaturan, perbarui aplikasi atau tonton video HD saat terhubung Wi-Fi.
  \item \textbf{Batang sinyal $\neq$ Kecepatan jaringan}: Sinyal kuat hanya berarti koneksi baik, kecepatan jaringan juga tergantung pada seberapa sibuk menara pemancar dan seberapa jauh Anda dari menara.
\end{itemize}

\textbf{Peringatan Penting: Matikan [Data Roaming] Sebelum Bepergian ke Luar Negeri!}\par

Menggunakan data seluler di luar negeri dapat menghasilkan tagihan mahal! Sebelum bepergian ke luar negeri, pastikan mematikan data roaming di pengaturan, atau beli kartu kuota lokal/paket roaming internasional terlebih dahulu.

\booksection{Wi-Fi: Koneksi, Berbagi, dan Keamanan}

Wi-Fi adalah andalan koneksi di rumah dan kantor, cepat dan tidak menghabiskan kuota. Namun Wi-Fi sebenarnya terus berevolusi, perbedaan antar standar cukup besar.

\booksubsection{Menghubungkan Wi-Fi:}

Buka [Pengaturan] > [WLAN] pilih jaringan masukkan kata sandi. Sekarang banyak ponsel mendukung koneksi dengan memindai kode, lebih nyaman.

\booksubsection{Berbagi Kata Sandi:}

Android dan iOS sama-sama memungkinkan ponsel yang sudah terhubung dengan cepat berbagi kata sandi ke teman. Android menghasilkan QR code, iOS bisa berbagi otomatis dengan mendekat.

\booksubsection{Evolusi Standar Wi-Fi:}

Standar Wi-Fi ditetapkan oleh IEEE, setiap generasi memiliki peningkatan signifikan:

\begin{center}
\begin{tabularx}{\linewidth}{|X|X|X|X|X|}
\hline
\textbf{Standar} & \textbf{Daya Maksimum} & \textbf{Pita Frekuensi} & \textbf{Kecepatan Maks} & \textbf{Karakteristik} \tabularnewline
\hline
Wi-Fi 4 & 802.11n & 2.4/5GHz & 600Mbps & Standar lama, cakupan lumayan \tabularnewline
Wi-Fi 5 & 802.11ac & 5GHz & 3.5Gbps & Saat ini paling umum, kecepatan cukup \tabularnewline
Wi-Fi 6 & 802.11ax & 2.4/5GHz & 9.6Gbps & Performa multi-perangkat kuat, latensi rendah \tabularnewline
Wi-Fi 6E & 802.11ax & 2.4/5/6GHz & 9.6Gbps & Menambah pita 6GHz, tidak padat \tabularnewline
Wi-Fi 7 & 802.11be & 2.4/5/6GHz & 46Gbps & Transmisi super cepat, agregasi multi-link \tabularnewline
\hline
\end{tabularx}
\end{center}

\begin{itemize}
  \item \textbf{Teknologi kunci Wi-Fi 6}: OFDMA (Orthogonal Frequency Division Multiple Access) memungkinkan beberapa perangkat mentransmisikan secara bersamaan, MU-MIMO (Multi-User Multiple Input Multiple Output) meningkatkan efisiensi konkurensi multi-perangkat, TWT (Target Wake Time) membuat perangkat lebih hemat daya.
  \item \textbf{Wi-Fi 6E}: Menambah pita frekuensi 6GHz pada basis Wi-Fi 6, menyediakan lebih banyak saluran, mengurangi interferensi. Membutuhkan dukungan ponsel dan router.
  \item \textbf{Wi-Fi 7}: Menggunakan MLO (Multi-Link Operation), dapat sekaligus memakai banyak pita frekuensi untuk transfer data, mengurangi latensi secara signifikan. Kanal super lebar 320MHz membuat kecepatan melonjak. Platform chip flagship seperti Snapdragon 8 Elite Gen5, Dimensity 9400 sudah mendukung Wi-Fi 7, ponsel seperti Xiaomi 17, Huawei Mate 80, Samsung S26 dapat memakai jaringan Wi-Fi 7 (perlu router Wi-Fi 7).
  \item \textbf{2.4GHz vs 5GHz}: 2.4GHz tembus dinding bagus tetapi lambat dan banyak interferensi; 5GHz cepat tetapi lemah menembus dinding. Pita 6GHz Wi-Fi 6E/7 paling cepat, tetapi paling lemah menembus dinding, cocok digunakan dalam ruangan yang sama.
\end{itemize}

\booksubsection{Risiko Wi-Fi Publik:}

Wi-Fi publik mungkin memiliki peretas yang menyadap data atau menyiapkan hotspot phishing. Aturan:

\begin{itemize}
  \item Hindari login ke rekening bank, memasukkan kata sandi.
  \item Sebaiknya gunakan VPN untuk koneksi terenkripsi.
  \item Jangan otomatis terhubung ke Wi-Fi asing.
  \item Gunakan situs web HTTPS (ada ikon gembok di address bar).
\end{itemize}

\booksection{Bluetooth: Serba Bisa Komunikasi Nirkabel Jarak Dekat}

Bluetooth mungkin salah satu teknologi nirkabel yang paling sering Anda gunakan di ponsel — headphone nirkabel, smartwatch, sistem mobil, keyboard dan mouse, semuanya terhubung melaluinya.

\booksubsection{Prinsip Teknis:}

Bluetooth bekerja di pita ISM 2.4GHz, menggunakan teknologi Frequency Hopping Spread Spectrum (FHSS) untuk menyebarkan data di 79 saluran dengan transmisi perpindahan cepat, secara efektif menghindari interferensi. Jarak komunikasi Bluetooth dibagi menjadi tiga kelas: Class 1 sekitar 100 meter, Class 2 sekitar 10 meter (umumnya digunakan ponsel), Class 3 sekitar 1 meter.

\booksubsection{Evolusi Versi:}

\begin{itemize}
  \item \textbf{Bluetooth 4.0}: Memperkenalkan Bluetooth Low Energy (BLE), memungkinkan perangkat baterai kecil seperti gelang kebugaran dan strap detak jantung juga terhubung nirkabel, daya tahan dari hari menjadi bulan.
  \item \textbf{Bluetooth 4.2}: Meningkatkan kecepatan transmisi dan perlindungan privasi, mendukung koneksi langsung IPv6.
  \item \textbf{Bluetooth 5.0}: Kecepatan transmisi dua kali lipat (2Mbps), jarak transmisi 4 kali (teoretis 300 meter), kemampuan siaran 8 kali. Mendukung dual audio (menghubungkan dua headphone secara bersamaan).
  \item \textbf{Bluetooth 5.1}: Memperkenalkan fungsi pengalamatan arah (AoA/AoD), akurasi posisi mencapai sentimeter.
  \item \textbf{Bluetooth 5.2}: Standar baru LE Audio, mendukung pengkodean LC3 (kualitas suara lebih baik, konsumsi daya lebih rendah), audio multi-stream, audio siaran.
  \item \textbf{Bluetooth 5.3/5.4}: Optimalisasi lebih lanjut pada konsumsi daya dan stabilitas koneksi, penyiaran periodik yang ditingkatkan.
\end{itemize}

\textbf{Adopsi LE Audio}: Pada 2025-2026, LE Audio telah memasuki tahap adopsi cepat. TWS baru merek mainstream (seperti Huawei FreeBuds Pro 5, Xiaomi Buds 6 Pro, Samsung Galaxy Buds3 Pro, dll.) sudah mendukung LE Audio dan codec LC3, memberikan latensi lebih rendah dan kualitas suara lebih baik. Saat membeli earphone baru, disarankan prioritaskan produk yang mendukung LE Audio.

\booksubsection{Headphone nirkabel}

\begin{itemize}
  \item \textbf{Transmisi audio}: Headphone nirkabel, audio mobil. LE Audio Bluetooth 5.2 membuat kualitas suara dan latensi melompat kualitas.
  \item \textbf{Transmisi data}: Mengirim file, foto jarak dekat. Meskipun kecepatan tidak sebaik Wi-Fi Direct, tetapi menang kesederhanaan.
  \item \textbf{Pasangan perangkat}: Smartwatch, gelang kebugaran, keyboard mouse, gamepad.
  \item \textbf{Layanan lokasi}: Bluetooth Beacon digunakan untuk navigasi indoor, pelacakan barang (seperti Apple AirTag).
  \item \textbf{Rumah pintar}: Jaringan Bluetooth Mesh mengontrol bohlam lampu, kunci pintu, dll.
\end{itemize}

\booksection{Remote Inframerah: Ponsel Berubah Jadi Remote Control Serba Bisa}

Remote inframerah adalah salah satu teknologi kontrol nirkabel tertua, tetapi masih praktis di ponsel. Sebagian ponsel Android (seperti model tertentu Xiaomi, Huawei) memiliki pemancar inframerah bawaan, dapat mengendalikan TV, AC, set-top box, dan perangkat rumah lainnya.

\booksubsection{Cara Kerja:}

Remote inframerah memancarkan sinyal pulsa cahaya inframerah frekuensi tertentu (biasanya 38kHz) melalui LED inframerah, penerima inframerah perangkat mendekode pulsa ini untuk menjalankan operasi yang sesuai. Merek berbeda menggunakan protokol pengkodean berbeda, yang umum adalah NEC, RC5, RC6, SONY, dll. Aplikasi inframerah di ponsel memiliki pustaka kode berbagai merek bawaan, pilih merek dan model untuk mencocokkan.

\booksubsection{Keterbatasan:}

\begin{itemize}
  \item \textbf{Directional}: Inframerah merambat lurus, perlu diarahkan ke perangkat, tidak bisa menembus dinding.
  \item \textbf{Jarak pendek}: Jarak efektif umumnya 5-8 meter.
  \item \textbf{Apple tidak mendukung}: iPhone tidak memiliki pemancar inframerah, tidak dapat menggunakan fitur ini.
\end{itemize}

\booksection{NFC: Keajaiban Sentuhan}

NFC (Near Field Communication, Komunikasi Jarak Dekat) adalah teknologi komunikasi nirkabel frekuensi tinggi dengan jarak sangat pendek (dalam 10 sentimeter), frekuensi kerja 13.56MHz, kecepatan transmisi tertinggi 424Kbps.

\booksubsection{Prinsip dan Mode Komunikasi:}

NFC berevolusi dari teknologi RFID, mendukung tiga mode kerja:

\begin{itemize}
  \item \textbf{Mode baca/tulis}: Ponsel membaca informasi dari tag NFC, seperti membaca stiker NFC untuk memicu aksi.
  \item \textbf{Mode peer-to-peer}: Dua perangkat NFC saling mentransmisikan data, seperti Android Beam (sudah dihentikan).
  \item \textbf{Mode kartu emulasi}: Ponsel mengemulasikan kartu pintar, ini mode yang paling sering digunakan — kartu transportasi, kartu akses, kartu bank semua menggunakan mode ini.
\end{itemize}

\booksubsection{Mekanisme Keamanan:}

\begin{itemize}
  \item \textbf{Jarak super pendek}: Jarak komunikasi dalam 10 sentimeter itu sendiri adalah jaminan keamanan fisik, penyadapan jarak jauh hampir tidak mungkin.
  \item \textbf{Chip keamanan}: Ponsel memiliki chip keamanan independen (SE) bawaan atau menggunakan HCE (Host Card Emulation), kunci pembayaran disimpan di area aman, bahkan jika ponsel di-root tidak dapat diekstrak.
  \item \textbf{Verifikasi sidik jari/wajah}: Saat pembayaran perlu konfirmasi biometrik, bahkan jika ponsel hilang tidak dapat disalahgunakan.
  \item \textbf{Tokenisasi}: Apple Pay, Huawei Pay, dll. tidak mentransmisikan nomor kartu asli, tetapi menggunakan token sekali pakai, bahkan jika dicegat tidak dapat digunakan ulang.
\end{itemize}

\booksubsection{Aplikasi Tipikal:}

\begin{itemize}
  \item \textbf{Kartu transportasi/kartu akses}: Menulis informasi kartu fisik ke ponsel, cukup sentuh ponsel.
  \item \textbf{Pembayaran mobile}: Apple Pay, Huawei Pay, dll., mengikat kartu bank ke ponsel, sentuh untuk membayar.
  \item \textbf{Pasangan perangkat}: NFC sentuh untuk memasangkan headphone Bluetooth, speaker dengan cepat.
  \item \textbf{Autentikasi identitas}: Pembacaan KTP NFC, verifikasi lisensi mengemudi elektronik.
\end{itemize}

\booksection{Hotspot Pribadi: Mengubah Ponsel Jadi Router Wi-Fi}

Saat kuota data Anda banyak, bisa mengaktifkan hotspot untuk tablet, komputer.

\begin{itemize}
  \item \textbf{Cara mengaktifkan}: Pengaturan > Hotspot Pribadi/Tethering, aktifkan.
  \item \textbf{Pengaturan dalam mode pengembang dapat mempengaruhi stabilitas sistem, opsi yang tidak dipahami jangan diubah sembarangan. Setelah mengaktifkan debugging USB, berhati-hatilah saat menghubungkan ke komputer, jangan hubungkan ke komputer yang tidak dipercaya.}: Sangat boros daya dan kuota, ingat untuk mengatur kata sandi yang kompleks, dan batasi jumlah perangkat yang terhubung.
  \item \textbf{USB tethering}: Menghubungkan komputer dengan kabel data, lebih stabil dan lebih cepat daripada hotspot nirkabel.
\end{itemize}

\begin{tcbpractice}{Praktikkan: Mengatur dan Mengelola Hotspot}
\begin{enumerate}
  \item Buka [Pengaturan], cari "Hotspot Pribadi" atau "Tethering", aktifkan, atur nama dan kata sandi.
  \item Gunakan perangkat lain (misalnya komputer) untuk mencari dan menghubungkan Wi-Fi ini.
  \item Di sisi ponsel, lihat apakah bisa mengelola perangkat yang terhubung, atau mengeluarkan perangkat yang tidak dikenal.
\end{enumerate}

\end{tcbpractice}
